%% kapa: exam=1 round=34 year=2023 period=1 form=A subject=real_estate subject_ko=부동산학원론 seq=11 number=91
%% source: Q-Net 감정평가사 기출문제 (cst00309) artlSeq=5211789 file=raw/1008/5211789/2023년도 제34회 감정평가사 1차시험 1교시.pdf
%% conversion: kordoc 4.24.51 -> tools/gen1.py (auto); review=proofread-sonnet; flags=none
\begin{question}{91}
부동산시장의 수요와 공급의 가격탄력성에 관한 설명으로 옳지 않은 것은? (단, 다른 조건은 동일함)
\begin{choices}
  \item 측정하는 기간이 길수록 수요의 탄력성은 더 탄력적이다.
  \item 공급의 탄력성은 생산요소를 쉽게 얻을 수 있는 상품일수록 더 탄력적이다.
  \item 수요의 탄력성이 탄력적일 경우 임대료가 상승하면 전체 임대수입은 감소한다.
  \item 대체재가 많을수록 수요의 탄력성은 더 탄력적이다.
  \item 제품의 가격이 가계소득에서 차지하는 비중이 작을수록 수요의 탄력성이 더 탄력적이다.
\end{choices}
\end{question}
\answer{5}
